% Section 7: Correctness Properties — Extended Formal Treatment
\section{Correctness Properties: Formal Guarantees}
\label{sec:rs-correctness-formal}

The Recursive Spawning Protocol rests on three interlocking correctness theorems that, taken together, establish an unconditional guarantee: that a spawn chain is an accountable, bounded, auditable, and finite lineage terminating exclusively in human judgment. These theorems are not heuristic. Each is a logical consequence of the structural constraints imposed by the three-layer BX3 architecture, enforced by mechanisms that no machine actor can bypass.

This section provides the full formal statements of the drift prevention, chain integrity, and termination theorems, complete proof sketches that expose the structural reasoning underlying each, and an analysis of the three-layer enforcement architecture that makes each proof valid.

% -------------------------------------------------------
\subsection{Drift Prevention: The Core Theorem}
\label{sec:rs-drift-prevention-formal}

\medskip
\noindent\textbf{Theorem 1 (Drift Prevention).} Let $C = \langle n_0, n_1, \ldots, n_k \rangle$ be a Recursive Spawning chain rooted at $n_0$ with human anchor $h(n_0) = h(n_k)$. Then for all $i$ ($0 \leq i \leq k$), the operating scope $R(n_i)$ is a descendant refinement of the intent of $h(n_0)$. Formally:

\[
\forall i \in [0, k]\ :\ R(n_i) \preceq_{h(n_0)} h(n_0)\mathrm{:intent}
\]

where $\preceq_{h(n_0)}$ denotes the descendant-refinement partial order induced by successive scope-subset operations. Equivalently: no node in $C$ can exhibit behavior outside the authorized intent of its human anchor.

\medskip
\noindent\textbf{Proof sketch.} The proof proceeds by structural induction on chain depth, where each inductive step is enforced at runtime by the Fact Layer pre-commit hook.

\medskip
\textit{Base case ($i=0$).} $R(n_0)$ is defined by $h(n_0)$ at provisioning. By the Purpose Layer's definition of authorized operating scope, $R(n_0)$ is a descendant refinement of $h(n_0)$'s intent by construction. No drift exists at the root.

\medskip
\textit{Inductive step.} Assume $R(n_i)$ is a descendant refinement of $h(n_0)$'s intent for some $i$ with $0 \leq i < k$. For $n_i$ to spawn $n_{i+1}$, the Fact Layer pre-commit hook verifies two conditions before the spawn is recorded:

\begin{enumerate}
    \item[(a)] $R(n_{i+1}) \subseteq R(n_i)$ --- the scope refinement constraint. The Purpose Layer's spawn authorization policy $P_{\mathrm{spawn}}$ requires strict scope subset; there is no authorized provision for lateral or upward scope expansion.
    \item[(b)] $\mathrm{depth}(n_{i+1}) \leq D_{\mathrm{max}}$ --- the depth constraint, also enforced by the pre-commit hook.
\end{enumerate}

The Fact Layer rejects any spawn operation violating (a) or (b). No actor---including the Bounds Engine---can execute a rejected spawn. Therefore $R(n_{i+1}) \subseteq R(n_i)$ holds for the child by structural enforcement, not by policy convention.

By the inductive hypothesis, $R(n_i)$ is a descendant refinement of $h(n_0)$'s intent. Since $R(n_{i+1}) \subseteq R(n_i)$, transitivity of subset gives that $R(n_{i+1})$ is also a descendant refinement of $h(n_0)$'s intent. Hence $n_{i+1}$ does not drift. By induction, no node in $C$ drifts. $\square$

\medskip
The drift prevention guarantee does not depend on the correctness or honesty of any machine actor. It depends on three structural conditions operating in conjunction:

\begin{enumerate}
    \item[(C1)] \textbf{Immutable parent pointer (Fact Layer enforcement).} $\alpha(n)$ is established once at provisioning and is thereafter immutable under the Fact Layer write protocol. Scope inheritance is traceable to a unique authorized lineage.
    \item[(C2)] \textbf{Forensic ledger (Fact Layer recording).} Every scope assignment is recorded in the append-only, hash-chained ledger. The ledger makes each step of the refinement chain empirically verifiable. The inductive hypothesis is not merely assumed---it is verified in the ledger at every step.
    \item[(C3)] \textbf{Exclusive human terminus (Purpose Layer enforcement).} The chain terminates at $h(n)$, which is non-migrating: $h(n_i) = h(n_0)$ for all $i$. There is no machine-actor alternative terminus. The Purpose Layer enforces $D_{\mathrm{max}}$ and $D_{\mathrm{ESCALATE}}$ as hard architectural bounds.
\end{enumerate}

If any of C1, C2, or C3 is absent, drift prevention fails. With all three present, drift is architecturally impossible. The proof sketch above shows exactly how C1 and C2 enforce the inductive step, and how C3 closes the inductive argument at a finite bound.

The contrapositive is also instructive: if drift is observed in any chain, at least one of C1, C2, or C3 has been violated. This makes the forensic ledger not merely an audit tool but a diagnostic instrument for identifying the specific structural failure that permitted drift.

% -------------------------------------------------------
\subsection{Chain Integrity: Tamper-Evident Topology}
\label{sec:rs-chain-integrity-formal}

\medskip
\noindent\textbf{Theorem 2 (Chain Integrity).} For any Recursive Spawning chain $C = \langle n_0, n_1, \ldots, n_k \rangle$, two properties hold simultaneously:

\begin{enumerate}
    \item[(I)] \textbf{Topological integrity:} $\alpha(n_i) = n_{i-1}$ for all $1 \leq i \leq k$.
    \item[(II)] \textbf{Node immutability:} No node configuration $(R(n_i), P_{\mathrm{spawn}}(n_i), h(n_i))$ is modified after provisioning.
\end{enumerate}

\medskip
\noindent\textbf{Proof sketch.} The proof addresses each property in turn.

\medskip
\textit{Topological integrity (I).} Each spawn event $n_{i-1} \xrightarrow{\mathrm{spawn}} n_i$ establishes $\alpha(n_i) = n_{i-1}$ as an immutable Fact Layer property at provisioning time. The Fact Layer write protocol rejects any operation attempting to modify $\alpha(n_i)$ after provisioning. Since the chain is constructed exclusively through authorized spawn events and no parent pointer is modified after establishment, $\alpha(n_i) = n_{i-1}$ holds for all $i \geq 1$ by construction. The parent pointer is not merely access-controlled---it is structurally immutable. There is no privileged actor, including the Purpose Layer after provisioning, capable of overwriting $\alpha(n)$. This eliminates the re-parenting drift pathway by making retroactive chain rewriting structurally impossible.

\medskip
\textit{Node immutability (II).} Each node $n_i$ is provisioned with a fixed configuration established by the Purpose Layer and recorded in the Fact Layer authorization ledger. The Fact Layer enforces immutability of these parameters after provisioning. The write protocol rejects any modification request from any actor---including the Bounds Engine and the Purpose Layer after provisioning. The Purpose Layer may issue a $\texttt{RESOLVE}$ directive that redirects a node's behavior, but this directive is recorded as a new ledger entry (a state transition), not as an alteration of the provisioned configuration.

The forensic ledger's hash-chaining provides evidentiary support for chain integrity: any attempt to insert a falsified spawn event or to reorder existing entries breaks hash continuity and is detectable by any observer with ledger read access. Chain integrity is not merely a property of the current chain state---it is a property of the entire historical record, enforceable at any future time. $\square$

\medskip
The chain integrity theorem has a practical consequence for auditors and incident investigators: the parent-pointer chain is not only intact but reconstructable at any future time from the forensic ledger. Given any node $v$, an auditor can reconstruct $C(v)$ by querying the ledger for all $\texttt{child-provisioned}$ events, extracting the $\alpha(\cdot)$ field from each entry, and traversing upward until reaching a node that does not appear as a child in any entry---that node is $h(v)$. The reconstruction is deterministic and reproducible; any two auditors who independently reconstruct $C(v)$ obtain identical results.

% -------------------------------------------------------
\subsection{Termination: Finite Bounded Escalation}
\label{sec:rs-termination-formal}

\medskip
\noindent\textbf{Theorem 3 (Termination).} Every Recursive Spawning chain $C = \langle n_0, n_1, \ldots, n_k \rangle$ terminates---either at a resolution state or at the escalation threshold---within a finite number of spawn steps bounded by $\max(D_{\mathrm{max}}, D_{\mathrm{ESCALATE}})$.

\medskip
\noindent\textbf{Proof sketch.} The proof establishes two lemmas and combines them.

\medskip
\textit{Lemma 1 (Depth-bounded growth).} The chain grows only through spawn events. By the Bounds Engine depth enforcement and the Fact Layer pre-commit hook, any spawn event producing $|C| \geq D_{\mathrm{max}}$ is rejected. Therefore the chain cannot exceed depth $D_{\mathrm{max}} - 1$ through autonomous spawn. Rejected spawns are logged; the Bounds Engine transitions to $\texttt{ESCALATING}$. Hence autonomous chain growth is bounded by $D_{\mathrm{max}}$.

\medskip
\textit{Lemma 2 (Escalation-triggered halt).} If the chain reaches $D_{\mathrm{ESCALATE}}$ before resolving the exception condition, the protocol enters $\texttt{ESCALATING}$ state and suspends all autonomous spawn activity. The Bounds Engine enters quiescent hold, awaiting a directive from $h(n_k)$---the human anchor. No further spawn events occur until $h(n_k)$ issues $\texttt{ACK}$, $\texttt{RESOLVE}$, or $\texttt{REJECT}$. Since $h(n_k) = h(n_0)$ is a human principal, the directive arrives in finite time. $\texttt{REJECT}$ terminates the chain definitively; $\texttt{RESOLVE}$ redirects behavior and records a new ledger entry; $\texttt{ACK}$ permits continued operation under continued Purpose Layer authorization.

\medskip
\textit{Theorem conclusion.} Combining Lemma 1 and Lemma 2: the chain grows autonomously at most $D_{\mathrm{max}} - 1$ spawn steps. If it reaches $D_{\mathrm{ESCALATE}}$ before resolution, it halts and escalates. If it reaches $D_{\mathrm{max}}$ before resolution, the Fact Layer rejects further spawn and triggers mandatory escalation. In all execution paths, the chain reaches a terminal state within $\max(D_{\mathrm{max}}, D_{\mathrm{ESCALATE}})$ spawn steps. Since both $D_{\mathrm{max}}$ and $D_{\mathrm{ESCALATE}}$ are finite integers defined at provisioning, the chain terminates in finite time. $\square$

\medskip
The termination guarantee precludes two failure modes that would defeat the upstream BX3 accountability guarantee:

\begin{enumerate}
    \item[(T1)] \textbf{Infinite spawn-escalate cycles.} When $D_{\mathrm{max}}$ is reached, the Fact Layer blocks further spawn and triggers mandatory escalation. The human anchor's $\texttt{REJECT}$ directive terminates the chain. There is no mechanism by which the chain can circumvent this bound and re-enter autonomous spawn---the Fact Layer pre-commit hook is structural, not configurable at runtime by any machine actor.
    \item[(T2)] \textbf{Escalation stalls.} The deterministic escalation path (Property~\ref{prop:deterministic-escalation}) means the escalation path is predictable and auditable before any exception occurs. The human anchor can determine the complete escalation path for any node by reading the forensic ledger entries along the chain. There are no runtime decisions that can alter the escalation path.
\end{enumerate}

\begin{property}[Deterministic Escalation Path]
\label{prop:deterministic-escalation}
For any node $v$ and any exception state $e$ at $v$, the Bailout Protocol's escalation path from $v$ to $h(v)$ is uniquely determined and is independent of the content of $e$ or the state of the Bounds Engine at any node on the path.
\end{property}

This property follows directly from parent-pointer immutability: the escalation algorithm traverses $\alpha(v), \alpha^{(2)}(v), \ldots, \alpha^{(k)}(v)$ where $k \leq D_{\mathrm{max}}$. Each $\alpha^{(i)}(v)$ is immutable and deterministic. The path does not depend on the exception type, severity, or any runtime state. It is a structural property of the spawn tree, not an operational decision made by any machine actor.

% -------------------------------------------------------
\subsection{Collective Guarantee: The Unified Accountability Theorem}
\label{sec:rs-unified-theorem}

Together, Theorems 1, 2, and 3 constitute the \textbf{Unified Accountability Theorem for Recursive Spawning}:

\medskip
\noindent\textbf{Theorem (Unified Accountability).} Let $C = \langle n_0, n_1, \ldots, n_k \rangle$ be a Recursive Spawning chain governed by the BX3 Framework. Then:

\begin{enumerate}
    \item $C$ is scope-faithful: for all $i$, $R(n_i)$ is a descendant refinement of $h(n_0)$'s intent (Drift Prevention, Theorem 1).
    \item $C$ is topologically stable: $\alpha(n_i) = n_{i-1}$ for all $i \geq 1$, and no node configuration is modified post-provisioning (Chain Integrity, Theorem 2).
    \item $C$ is finite: it terminates within $\max(D_{\mathrm{max}}, D_{\mathrm{ESCALATE}})$ steps at a human principal (Termination, Theorem 3).
    \item $C$ is tamper-evident: the complete history of all spawn events and state transitions is recorded in the append-only forensic ledger with cryptographic hash-chaining (Fact Layer).
\end{enumerate}

These four properties are mutually reinforcing. Scope-faithfulness without topological stability is meaningless: an actor could drift by re-parenting. Topological stability without scope-faithfulness is meaningless: a chain could be stable but operating outside authorized scope. Finiteness without scope-faithfulness and topological stability is meaningless: a finite chain could still be operating outside authorized scope. The forensic ledger is the evidentiary foundation that makes all three properties independently verifiable.

The Unified Accountability Theorem establishes that a Recursive Spawning chain governed by the BX3 Framework is a provably accountable, integrity-preserving, and finite lineage of agents---suitable for deployment in high-stakes operational contexts where auditability and correctness are non-negotiable requirements.